\section*{Chapter \ref{ch:hf:the-business-of-market-making} --- The Business of Market Making}

\begin{solution}{exo:hf:the-business-of-market-making:1}
$2{,}145.3/248.5=\$8.63$ million a trading day.
\end{solution}

\begin{solution}{exo:hf:the-business-of-market-making:2}
$10^4\times485.8\times10^6/(1{,}940\times10^9)=2.50$ basis points. The numerator includes trading income from other products, so the ratio is an
upper bound on what the firm earned per euro of exchange-traded products.
\end{solution}

\begin{solution}{exo:hf:the-business-of-market-making:3}
$V^\ast=40{,}000/(0.3\times10^{-4})=\$1.33$ billion a day.
\end{solution}

\begin{solution}{exo:hf:the-business-of-market-making:4}
$\pi=0.3\times10^{-4}\times2\times10^9-40{,}000=\$20{,}000$; leverage $V/(V-V^\ast)=2/(2-1.333)=3.0$. At \$1.8 billion, $\pi=54{,}000-40{,}000=
\$14{,}000$: 30\% less profit for 10\% less volume.
\end{solution}

\begin{solution}{exo:hf:the-business-of-market-making:5}
The mid at the moment of a fill is stale: fills happen disproportionately when the efficient price has already moved through the quote (informed
traders act on the gap, and quotes are slow to follow). Against the mid the Quoter always earns half the spread at the fill; against the efficient
price it has already lost 0.173 cents, and the mid catches up over the next seconds, which is the adverse selection measured against the mid.
\end{solution}

\begin{solution}{exo:hf:the-business-of-market-making:6}
Between 10 seconds ($+0.024$ cents) and 20 seconds ($-0.060$). Any position held longer than about ten seconds has, on average, lost its half-spread:
the Quoter must earn the next spread, or hedge, before then.
\end{solution}

\begin{solution}{exo:hf:the-business-of-market-making:7}
$-0.150$ cents a share: the $-0.050$ with the rebate, less the 0.100 cent rebate. Fees do not enter the Quoter's decisions, so orders, fills and
prices are identical; only the fee line changes.
\end{solution}

\begin{solution}{exo:hf:the-business-of-market-making:8}
One hour is one draw. The hourly P\&L of this kind of quoter has a standard deviation near \$89, so \$61 is well inside the noise of a quoter whose
true mean is negative (the chapter's is $-\$19$). Decompose the P\&L, run many independent hours, and report the mean with its standard error
before sizing anything.
\end{solution}

\begin{solution}{pb:hf:the-business-of-market-making:1}
\begin{enumerate}
\item A firm providing liquidity by algorithm across many instruments and venues; it sells immediacy, priced at half the spread.
\item \$8.63 million a day; market making \$1,666 million, 77.7\% (\$6.71 million a day).
\item Staff \$528.1, communication and data \$249.2, operations \$97.9, depreciation \$64.4: \$939.6 million, \$3.78 million a trading day.
\item \$2.09 million of adjusted net trading income and \$0.51 million of pay per employee.
\item 2.50 basis points; \euro0.32 million of fixed costs per employee; EBITDA margin 41\%.
\item Trading income net of volume-driven fees per unit traded; $\pi(V)=(c-v)V-F$.
\item $V^\ast=F/(c-v)$; $d\ln\pi/d\ln V=V/(V-V^\ast)$.
\item \$1.33 billion; \$20,000 a day and a leverage of 3.0 at \$2 billion.
\item $250\times20{,}000=\$5.0$ million.
\item $50{,}000/(0.3\times10^{-4})=\$1.67$ billion.
\item $40{,}000/(0.2\times10^{-4})=\$2.0$ billion; profit zero at \$2 billion.
\item See \cref{prop:hf:the-business-of-market-making:split}: the three brackets telescope to $m(T)-p_i$.
\item Its own orders move the top of the book; joining that top chases its own orders into a message storm (336,012 messages against 323 in the
tutorial's run).
\item 38,270 shares an hour, 11\% of the market's volume, 1,256 messages.
\item Spread $+0.487$, adverse selection $-0.463$, inventory $-0.174$, fees $+0.100$: $-0.050$ cents a share.
\item Its fills cluster when the efficient price has moved through its quote; the mid at the fill is stale.
\item Standard deviation \$89 an hour, standard error $89/\sqrt{10}=\$28$ for a mean of $-\$19$: not distinguishable from zero.
\item The spread pays and the position costs (\euro1.55 against \euro0.68 per trade); speed rank drives performance. This Quoter earns its spread and
loses it on positions held too long, and has no way to be faster than the informed flow.
\item Break-even at \$1.33 billion a day; operating leverage 3.0 at \$2 billion a day.
\item Lower fixed costs and a lower break-even volume first: capture is set by competition and volume by the market, while a small margin over a
large fixed base makes every quiet month a loss.
\end{enumerate}
\end{solution}

\begin{solution}{iq:hf:the-business-of-market-making:1}
Immediacy: the ability to trade now. The investors who cross the spread pay for it, half the spread per share; the market maker's cost is the
adverse selection from the ones who know more.
\iqlookfor{immediacy, half-spread, and that informed flow is the cost.}
\end{solution}

\begin{solution}{iq:hf:the-business-of-market-making:2}
Spread capture $s\,q\,(m(t)-p)$, adverse selection $s\,q\,(m(t+H)-m(t))$, inventory $s\,q\,(m(T)-m(t+H))$. The horizon separates the price move
that the fill predicted (information) from later moves of the held position (risk); the choice of $H$ moves P\&L between the two terms, never the
total.
\iqlookfor{the telescoping identity and that $H$ is a convention.}
\end{solution}

\begin{solution}{iq:hf:the-business-of-market-making:3}
Adverse selection: the fills come from traders who know the next price. Try a better fair price than the mid, skew or widen when the book or the
lead instrument signals a move, and hedge or unwind within the horizon over which the mark-out decays.
\iqlookfor{mark-out curves, a fair price, and acting within the decay horizon.}
\end{solution}

\begin{solution}{iq:hf:the-business-of-market-making:4}
A feedback loop: the quoter reads a book that contains its own orders and keeps repricing to join or leave itself. Quote off the book without one's
own orders, and add a minimum price change before repricing.
\iqlookfor{self-reference in the market-data view, and a message-rate test.}
\end{solution}

\begin{solution}{iq:hf:the-business-of-market-making:5}
Fixed costs: profit is $(c-v)V-F$, so its percentage change is $V/(V-V^\ast)$ times that of volume; and capture itself varies with volatility and
competition.
\iqlookfor{operating leverage and the break-even volume.}
\end{solution}

\begin{solution}{iq:hf:the-business-of-market-making:6}
$2\times89/\sqrt n=19$ gives $n=(178/19)^2\approx88$ hours.
\iqlookfor{standard error scaling and a number.}
\end{solution}
